% Source-aligned English for printed p.157.
\SourcePage{157}

\LayerTarget{tohoku:def:sections-and-supports}
\noindent sheaf and $A$ is a subset of $X$, write $\Gamma(A,F)$ for the group
of sections of $F$, regarded as an étalé space, over $A$, and put
$\Gamma(F)=\Gamma(X,F)$. Thus $\Gamma(A,F)=\Gamma(F\vert A)$, where
$F\vert A$ denotes the restriction of $F$ to $A$. More generally, let $\Phi$
be a nonempty upward-filtered set of closed subsets of $X$ such that
$A\in\Phi$ and $B\subset A$ imply $B\in\Phi$. For short, we call $\Phi$ an
\textit{antifilter of closed subsets} of $X$. Let $\Gamma_\Phi(F)$ be the
subgroup of $\Gamma(F)$ formed by the sections $f$ whose support---the
complement of the largest open subset of $X$ on which $f$ restricts to
zero---belongs to $\Phi$.
\LayerTarget{tohoku:def:cohomology-with-supports}
The functor $\Gamma_\Phi$ from $\mathbf C^X$ to abelian groups is visibly left exact.
We must therefore consider its right derived functors, equivalently its right
satellites; see 2.3. They exist by the corollary to Proposition 3.1.1. Denote
them by $H^p_\Phi(X,F)$, where $p$ is any integer. The theory of derived
functors gives:

\LayerTarget{tohoku:prop:3.2.1}
\indent\textsc{Proposition 3.2.1.} \textit{The system of functors
$H^p_\Phi(X,F)$ \textup{(}$-\infty<p<+\infty$\textup{)} is characterized by the following
conditions: it is a cohomological functor on $\mathbf C^X$ with values in abelian
groups; $H^p_\Phi(X,F)$ is zero for $p<0$, agrees with $\Gamma_\Phi(F)$ for
$p=0$, and is effaceable---that is, it vanishes on injective $F$---for
$p>0$.}

\LayerTarget{tohoku:construction:compute-cohomology-with-supports}
To compute $H^p_\Phi(X,F)$, take a right resolution of $F$ by injective
sheaves $C^i$, or more generally by sheaves known in advance to satisfy
$H^p_\Phi(X,C^i)=0$ for $p>0$; such $C^i$ are called
$\Gamma_\Phi$-acyclic. If $C(F)$ denotes the resulting complex, then
\[
H^p_\Phi(X,F)=H^p\Gamma_\Phi(C(F)).
\]
To use this method, it is important to have criteria ensuring that a sheaf is
$\Gamma_\Phi$-acyclic; some are given below in 3.3.

\LayerTarget{tohoku:construction:inverse-image-comparison:p1}
Let $f:X\to Y$ be a continuous map. For every sheaf $F$ on $Y$,
one naturally defines the \textit{inverse-image sheaf} of $F$ under $f$,
written $f^{-1}(F)$ by abuse of notation: regarding $F$ as an étalé space
over $Y$, one uses the ordinary definition of the inverse image of a fibered
space. This gives an \textit{exact additive covariant functor}
$F\to f^{-1}(F)$ from $\mathbf C^Y$ to $\mathbf C^X$, which we denote by $\mathrm{g}$. For
every $F\in \mathbf C^Y$, there is an evident natural homomorphism
$\Gamma(F)\to\Gamma(\mathrm{g}(F))$; equivalently, there is a natural
homomorphism of functors $\Gamma^Y\to\Gamma^X\mathrm{g}$, where the
superscript records the space relative to which $\Gamma$ is taken.
More generally, let $\Phi$ be an antifilter of closed subsets of $X$ and
$\Psi$ an antifilter of closed subsets of $Y$ such that $B\in\Psi$ implies
$f^{-1}(B)\in\Phi$. Then for every $F\in \mathbf C^Y$, the homomorphism
$\Gamma(F)\to\Gamma(\mathrm{g}(F))$ sends $\Gamma_\Psi(F)$ into
$\Gamma_\Phi(\mathrm{g}(F))$, giving a natural homomorphism
$\Gamma_\Psi\to\Gamma_\Phi \mathrm{g}$. Since $\mathrm{g}$ is \textit{exact},
$(R\Gamma_\Phi)\mathrm{g}$ may be regarded as a \textit{cohomological functor} on $\mathbf C^Y$ whose
degree-zero part is $\Gamma^X_\Phi \mathrm{g}$. Since the cohomological functor
$R\Gamma^Y_\Psi$ is ``universal'' by Proposition 2.2.1, the natural
homomorphism $\Gamma^Y_\Psi\to\Gamma^X_\Phi \mathrm{g}$

\TohokuContinuation{158}

